\section{Who can bind the state that enforces the rule?}
\label{sec:16.7}
A state can bind a private developer because it controls something the developer needs. When the state is the operator, the institution that writes the rule, controls the evidence and supplies the remedy is the institution the rule is meant to constrain.

Standing to sue has never turned on citizenship, and the doctrines closing American courts to targeting claims — political question, collapsed damages remedies, state secrets, statutory exceptions for injuries abroad and combatant activities — are indifferent to who the plaintiff is. They close the forum for a subject matter, and they've closed it to citizens. Three Americans killed in Yemen in 2011, one a sixteen-year-old nobody alleges was targeted: the court held that a plausible due-process claim had been stated as to the one who was targeted, and that no remedy was available for it \autocite{alaulaqi2014panetta}. A separate suit by an American journalist who survived five drone strikes was dismissed for want of standing to sue after a state-secrets dismissal below \autocite{kareem2021}.

And those cases concentrate in one appellate circuit by venue and by where the defendants sit — so the ordinary correction mechanism, several circuits disagreeing and forcing a higher court's attention, doesn't run. The doctrine hardens without ever being tested against a competing reading.

The mechanisms can work as mechanisms and fail as checks. In the Israeli case the apparatus performed — journalists published, serving personnel spoke, the courts were trying a sitting prime minister. No finding attached to a body with standing to stop what it described.

The American version: a civilian-harm investigation into Minab was, by officials' account, effectively complete for months and not released, while a Pentagon spokesperson said it remained ongoing \autocite{Bloomberg2026KillChain,fpj2026minabinquiry} — and the lever available to the Senate committee was withholding three-quarters of the Defense Secretary's travel budget until the unredacted investigation was handed over \autocite{stassis2026senate}.

What the Pentagon wouldn't publish, the UN mission did. The finding is public, reasoned, and attached to nothing — no jurisdiction, no answer to the mission's requests for information \autocite[para.~97]{UNFFMIran2026} and none to its findings beyond a statement that engaged none of the evidence \autocite{rferl2026whitehouseun}, the domestic investigation still unreleased.

The courts show the same pattern: transparency litigation succeeded where damages litigation failed. A court of appeals compelled disclosure of the executive's own legal analysis of targeted killing \autocite{nytimes2014olc}. The judiciary will make a government explain the legal basis on which it killed somebody and will not decide whether the killing was lawful.

What matters here is rules binding a state toward people who'll never vote in its elections, the purpose of international humanitarian law; what it hasn't supplied is a forum, and the bin Ali Jabers' case shows the forum switched off at home as well, by a decision that never had to reach who was suing.

All but one of the forum closures above are Congress's to change. A statute can create a cause of action for people harmed by a deployment that crossed its floor; waive sovereign immunity for it, including the Federal Tort Claims Act's exceptions for combatant activities and for claims arising abroad \autocite{usc2680}; and replace dismissal on state-secrets grounds with the procedure the Classified Information Procedures Act already uses in criminal cases, in which a judge reviews the evidence, the case goes on, and a government that still withholds evidence bears the cost \autocite{cipa1980}. None of that needs a new court, and a new court without it would reproduce \emph{al-Aulaqi}: a plausible claim and no remedy. Under capture, a refusal's remaining work is a record, and a record matters only if it can reach whoever made the decision. Under \emph{Trump v. United States} \autocite{trump2024immunity}, the President is absolutely immune for acts within his core constitutional powers and at least presumptively immune for his other official acts, commanding the armed forces among them, and testimony or records probing immune conduct can't be put before a jury. So the record reaches every commander except the one at the top, and a statute saying otherwise would be reviewed by the Court that made the holding. Making the record reach him means changing the Court. Appendix~\ref{sec:B} sets out a program for that.

The chapter's cases are American, and “the state” is tested here on one constitutional order. Europe is a different one: the Convention applies, has extraterritorial reach, and produces judgments states sometimes obey; there is no state-secrets privilege of the American kind and no collapsed damages remedy. The AI Act makes the boundary visible: it can require documentation from private providers, subject high-risk systems to review, and impose penalties tied to global turnover — and it excludes military, defence and national-security uses from its scope. Public authorities stay covered elsewhere, but enforcement against a member state's agencies still runs through institutions that state constitutes.

For review to bind here, the chain that makes review binding has to run against a state: initiation by people outside its electorate, evidence it would withhold, an interim order before the use is complete, and a remedy that doesn't need its consent. A principle supplying none of those may state the rule correctly and identifies nobody who can make the state obey it.

A treaty obligation of that kind exists: the Council of Europe's Framework Convention, binding at the level of a human-rights instrument with operational detail left to domestic law \autocite{coe2024framework} — and not yet in force, needing five ratifications, of which the EU deposited one in May 2026 \autocite{coe2026euratifies}. Everything else at that level is voluntary.

Every instrument I propose binds through one of three channels: a licence term, an appropriation condition, or a statute. The first two run through procurement, and procurement is one office. A captured Congress waives the appropriation condition; one memorandum from the Office of Management and Budget strikes the licence terms from every agreement at once. Three second routes exist and none is as strong. A collective agreement carries codetermination and the individual refusal right without the state. A coalition partner's precondition carries the review ratio into any war the partner takes part in. State and municipal procurement carry the civilian terms into deployments the federal government doesn't run. The force term has no second route, which is why it is a statute.

\textbf{Residual risk.} One residual risk no design choice removes, and it's what the floor was introduced for. Every lever belongs to somebody, and a government that can reach whoever holds the lever has the lever. A constraint that holds against the one in possession has to be held by something it cannot instruct, and under total capture no instrument is. Every instrument in this book ends at this limit. People don't end there. Under the most complete terror of the last century most people complied, some did not, and some were left to tell the story.
